% Expanded replacement for Part IV only.
% Replace the former Part IV in its entirety with this section.
% Rename the image files in \includegraphics{} to match the figures folder
% used by the manuscript.

\section{Technical Approach}

\subsection{A. Scenario-Scoped, Terrain-Aware Simulation-to-Omniverse Integration}

Within ECLIPSE, NVIDIA Omniverse is the environment in which a completed mission scenario becomes an inspectable three-dimensional operational scene. Its role is not simply to render CAD models. To reproduce the scenario created upstream, the visualization must receive the selected module geometry, the position and orientation of each placed instance, the lunar terrain associated with the site, rover routes, and the time-indexed operational events generated by the DES. These inputs originate in several ECLIPSE products: the SysML-derived system representation, Urban Planning outputs, USD scene assets, waypoint files, and DES results. The central integration challenge is therefore to convert these distributed products into one unambiguous visualization definition that can be loaded without manually rebuilding the scene.

Before this work, the Omniverse playback demonstrated that ECLIPSE could visualize an ISRU mission, but it was not a reusable integration contract. The extension could display a prepared scenario, yet it did not receive a formal, scenario-scoped description of the mission from the upstream pipeline. Module placement, rover timing, route selection, and the simulation data used for playback could consequently depend on manually prepared visualization logic. A change to the mission architecture therefore risked creating a mismatch between the scenario evaluated by the DES and the scenario displayed in Omniverse.

\begin{figure}[htbp]
    \centering
    \includegraphics[width=0.80\linewidth]{figures/omniverse_old_architecture}
    \caption{Previous Omniverse playback architecture. Scene assets, waypoint routes, and DES outputs were combined through a visualization scenario prepared specifically for the ISRU proof of concept.}
    \label{fig:omniverse-old-architecture}
\end{figure}

The revised workflow introduces a visualization manifest as the explicit interface between the ECLIPSE pipeline and Omniverse. After a DES run, the workflow generates a scenario package containing the scene USD, waypoint USD, terrain reference, system representation, DES result package, time-indexed DES log, actor definitions, route information, and terrain-projection diagnostics. These artifacts are written under the named scenario output directory rather than to one global location. Multiple completed scenarios can therefore coexist, and a subsequent run does not overwrite the scene, routes, or telemetry used to visualize a previous scenario. The manifest identifies the exact files that belong together and records how they should be interpreted by the extension.

\begin{figure}[htbp]
    \centering
    \includegraphics[width=0.80\linewidth]{figures/omniverse_manifest_architecture}
    \caption{Scenario-scoped manifest architecture. The generated manifest identifies the scene, terrain, routes, actors, and DES outputs consumed by the Omniverse extension.}
    \label{fig:omniverse-manifest-architecture}
\end{figure}

The extension acts primarily as a package reader and playback client. It lists the available scenario manifests, loads the package selected by the user, creates or updates actors from the manifest, and plays their state against the time-indexed DES log. It provides route visibility controls and overview, follow, and rover-chase camera modes. The extension also displays mission and rover telemetry corresponding to the current playback time. This dashboard is time-synchronized with recorded DES results; it is not a direct stream from a DES running concurrently. This distinction is important because it preserves a deterministic connection between a completed scenario result and its visual playback. Some compatibility logic remains for the established ISRU telemetry schema, so the extension is manifest-driven rather than entirely free of scenario-specific assumptions.

Spatial coherence requires more than referencing the same files. Urban Planning provides a metric map frame for module placement and route generation. This map frame is propagated into the visualization package so that the prepared terrain asset, module coordinates, and route coordinates use a common reference. The manifest builder aligns the terrain USD to that frame and constructs a terrain sampler from the mesh. The sampler identifies the terrain triangle beneath a projected point and interpolates the local terrain elevation. It is used both to ground static modules and to lift rover routes onto the surface.

For static assets, grounding is based on more than the center point of the model. Points distributed across the scaled CAD footprint are sampled on the terrain. A least-squares plane is fitted to these local terrain samples, producing the average slope under the module. The module receives the corresponding pitch and roll, then a final vertical adjustment ensures that its transformed footprint does not penetrate the terrain. The process uses the resolved CAD bounding box, including internal geometry transforms, rather than assuming that the origin of every source model is located at its physical base. This removed several earlier placement failures in which a model appeared below the terrain or floated above it because its internal geometry was offset.

Routes are treated similarly. The planned route between upstream waypoints is densified, projected to the terrain mesh, and evaluated segment by segment. The extension can render slope-classified segments and waypoint markers, allowing the user to see whether the terrain used for visualization is consistent with the route-feasibility assumptions used during Urban Planning. Rover yaw and pitch are calculated from the local tangent of the projected path, and motion is interpolated by distance along the route rather than by the index of a waypoint. The DES uses the scenario route length and selected route-level metrics; it does not yet execute a waypoint-by-waypoint vehicle dynamics model. Likewise, the implemented workflow aligns the terrain asset supplied for a scenario, but automatic retrieval and clipping of arbitrary terrain patches from a global lunar terrain database remains future work.

\begin{figure}[htbp]
    \centering
    \includegraphics[width=0.92\linewidth]{figures/terrain_route_slope}
    \caption{Terrain-projected module placement and slope-classified rover route. The terrain sampler grounds module footprints and projects route geometry to the prepared lunar mesh.}
    \label{fig:terrain-route-slope}
\end{figure}

\subsection{B. Multi-Format CAD Ingestion and SysML Metadata Update}

CAD geometry is the visual representation of an architectural element in ECLIPSE, but it can also contain information useful for engineering analysis. Before this work, the Model Submission step expected a USD-family file that could be referenced directly by the scene builder and Omniverse. The interface could inspect visualization-oriented stage metadata, such as resolved mesh bounds, units, materials when available, and the USD up axis. This approach worked for prepared USD assets already compatible with Omniverse, but it required a stakeholder who owned a STEP, STL, or OBJ model to convert it outside ECLIPSE before submission. It also did not establish a consistent mechanism for distinguishing metadata from the original engineering file from metadata of the visualization asset, or for returning newly extracted physical information to the SysML authoritative source of truth.

The revised CAD submission workflow addresses both limitations. Its first objective is to make stakeholder geometry usable by the pipeline even when it was not delivered in a USD-family format. The implemented workflow accepts STEP/STP, STL, OBJ, and USD-family files. STEP/STP inputs are read through OpenCascade using pythonOCC and tessellated into meshes. STL and OBJ are read as meshes through trimesh. USD-family files are opened through the USD Python API. Each supported submission is standardized into an Omniverse-compatible USD asset and a GLB preview for the browser interface. The workflow is extensible through additional format readers, but it does not yet support every possible CAD format.

Conversion alone is insufficient because engineering information can be absent from, changed by, or difficult to recover after a mesh conversion. The workflow therefore preserves two complementary metadata products. Source-CAD metadata describes the submitted artifact, including source name, file format, units when available, surface area, center of mass, and volume only when the geometry is sufficiently watertight to support a reliable calculation. Visualization metadata is generated after the USD asset is available and records the information needed to construct and inspect the scene: generated USD and preview paths, resolved bounding box, stage units, mesh count, available materials, and orientation information. For converted source formats that do not provide a dependable orientation convention, the browser preview allows the user to specify the up axis and front axis. The front-axis selection defines the forward direction needed for moving actors, while the up axis determines the source orientation used during conversion and placement.

\begin{figure}[htbp]
    \centering
    \includegraphics[width=0.95\linewidth]{figures/cad_metadata_pipeline}
    \caption{Multi-format CAD conversion and dual metadata propagation. Source-CAD metadata and post-conversion visualization metadata are kept distinct, while scaled physical metadata is written to the SysML authoritative source of truth.}
    \label{fig:cad-metadata-pipeline}
\end{figure}

SysML continues to govern the intended architectural dimensions of each module. CAD bounding boxes therefore do not overwrite SysML dimensions. Instead, the CAD geometry is scaled to fit the SysML target dimensions. The scale is applied consistently to scale-dependent physical metadata: surface area scales with the square of the length factor, volume with the cube, and center-of-mass coordinates with the transformed geometry. The resulting values, together with provenance such as the source file and format, are written into an auto-generated CAD metadata block in the corresponding SysML part and mirrored in the associated JSON asset. This approach preserves the distinction between the mission architecture, which is defined by SysML, and the geometry supplied to visualize it.

The standardized USD representation is subsequently referenced by the scene and manifest builders. The resolved dimensions and orientation metadata are used by the terrain-grounding procedure, while the front-axis selection allows rover models from heterogeneous sources to align with the direction of travel. The interface can also reuse CAD files already stored in the repository CAD library. Integration with a separate curated engineering and lunar metadata library is a planned extension, not a completed capability.

\subsection{C. DES: From an Opaque ISRU Demonstration to a Configurable and Extensible Framework}

Within ECLIPSE, the DES is the operational layer that converts a static mission architecture into time-dependent mission behavior. Starting from the active systems, their interfaces, their placement, and the routes created during Urban Planning, it determines when resources are produced, transported, stored, transformed, or consumed. It also evaluates power demand and availability and produces the time histories and measures of effectiveness used to assess a scenario. These outputs support both quantitative comparison in the web interface and temporal playback in Omniverse. For the DES to serve as a genuine mission-design tool rather than only an analysis script, its operating assumptions must be visible, configurable, and traceable to the scenario defined by the user.

The original DES was appropriate as a proof of concept but unsuitable as a general mission-definition interface. It represented a predetermined ISRU chain: regolith excavation and hauling, ISRU processing, LOX accumulation, LOX transport, and selected power consumers. Several operating choices were embedded in Python, including rover roles, storage behavior, thresholds, timing, and power events. The user could vary a small set of sliders, but could not inspect the full operational logic or connect a selected architecture to another resource flow. The scenario therefore behaved as a black box: an architecture was provided to the DES, but the detailed assumptions that converted it into operations remained mostly invisible.

The first and foundational step was to replace the ISRU DES black box with an explicit scenario definition that a user can inspect and configure. Previously, selecting an ISRU architecture mainly activated a fixed Python controller. In the revised workflow, the user first defines the operational architecture itself: placed module instances and multiplicities, SysML interfaces, resource and power connections, rover assignments, and rover routes. The interface makes these relationships visible on the mission map and in the scenario-logic view. It exposes the equations associated with each module and guides the user toward the equations required for its operational role. The objective is not merely to add more sliders, but to make the assumptions that transform a set of selected modules into an operational ISRU mission intelligible and traceable before a simulation is run.

The numerical configuration then exposes the ISRU assumptions that were previously embedded in code: simulation duration; rover fleet counts, payloads, and route assignments; transport and processing coefficients; plant and storage settings; power generation; battery state; and module power-demand models. A user can represent module demand as an average value, a time profile, or a restricted equation. All equations and numerical parameters governing the ISRU scenario are exposed through the interface: they are either defined in the scenario-logic and equation views or edited in the DES configuration view. No mission-specific ISRU assumption must remain hidden in the Python controller for the user to modify the scenario; Python retains the execution mechanisms that apply the selected configuration. Scenario validation checks the consistency of the placed architecture, required interfaces, assignments, and routes before execution, and reports missing or incompatible elements. These user-visible choices are assembled into a structured scenario configuration passed to the DES workflow and preserved with the resulting named scenario. A completed configuration can consequently be reloaded, inspected, modified, and compared to another completed configuration without reconstructing its operational assumptions manually. Within the scope of the ISRU behavior engine, the simulation is therefore no longer an opaque fixed scenario: its architecture, routes, equations, numerical parameters, and power assumptions are accessible to the end user.

\begin{figure}[htbp]
    \centering
    \includegraphics[width=0.98\linewidth]{figures/isru_scenario_configuration}
    \caption{User-visible ISRU scenario definition and DES configuration. The interface exposes asset multiplicity, rover routes, SysML interfaces, module equations, and numerical simulation assumptions that were previously hidden in the proof-of-concept model.}
    \label{fig:isru-scenario-configuration}
\end{figure}

The ISRU engine itself was extended rather than discarded. It supports multiple regolith rovers, multiple LOX rovers, and multiple ISRU plant instances. Each plant has separate regolith input storage, and available regolith rovers are dispatched toward the least-filled eligible plant. The selected scenario routes provide the effective haul distances used by the DES, while rover cycles include outbound travel and empty return travel. The LOX route and pickup threshold determine when an available LOX rover is dispatched. The model represents routes through their effective lengths and selected route metrics, rather than through detailed vehicle dynamics at each terrain waypoint.

Power assumptions were also made explicit. Solar generation, battery capacity, initial charge, management interval, and module demand models can be configured. ISRU processing and LOX-storage energy are now recorded as pending demands and passed to the power manager, addressing the earlier situation in which major ISRU energy use could be reported without contributing to the power-feasibility calculation. The implemented power model remains deterministic and simplified: it does not yet represent a detailed electrical network, equipment failures, or complete rover-charging operations.

The second step was to generalize the mission-description language beyond the ISRU template. The purpose was not to create a separate Python controller for every possible lunar mission. Instead, the workflow treats a submitted SysML architecture as a graph whose nodes are modules and whose directed interfaces represent resource and power relationships. From this graph, the system identifies a limited set of generic operational roles. It does not need prior semantic knowledge of the name ``Propellant Depot,'' for example. If the same resource is received, retained, and later sent onward, the graph and its interfaces identify a Storage behavior. Similarly, a module with an outgoing resource but no corresponding incoming resource is interpreted as a Source, while a module that receives one resource, produces another, and requires a processing time is interpreted as a Processor.

Step 4 serializes the placed module instances, SysML interfaces, resource routes, rover assignments, module classes, equations, and power configuration into a scenario description. Classes are inferred from graph connectivity, port direction, resource-flow direction, and relevant physical attributes; they are not intended to be inferred only from module names. Because a graph can remain ambiguous, the interface displays the inferred class for every module and allows the end user to override it. This creates a controlled collaboration between automated interpretation and engineering judgment: the pipeline provides a consistent initial behavior model, while the user remains responsible for correcting an interpretation that does not reflect the intended mission.

Each role corresponds to a standard SimPy behavior and a defined set of required relationships. A Source makes an outgoing resource available when it is requested. A Processor waits for its required input, applies the configured resource transformation, waits for the associated processing time, and places the resulting output into its output inventory. A Storage receives and retains a resource, then makes the same resource available to downstream connections. A Consumer receives power or resources without transforming them. A Transporter represents a rover fleet that carries a selected resource between two connected modules along an assigned route, including its return travel. A PowerGenerator supplies the power-accounting model. The role therefore gives the compiler enough operational structure to create processes, inventories, queues, transport resources, and power demands without requiring a mission-specific controller.

The role also tells the interface which equations are meaningful. Rather than asking the user to write unrestricted code, the equation builder guides the construction of resource-output laws, processing-time laws, storage relationships, travel-time laws, and energy or power laws using ports, attributes, and scenario variables exposed as selectable tokens. The scenario definition is checked before execution: the validator verifies that resource connections have compatible sources and destinations, transport links have assigned routes, required role equations are present, and equation outputs correspond to implemented simulation effects. In this way, the user is not only editing a numerical input file; the user is defining a mission behavior in a constrained, interpretable language that can be compiled into DES processes.

The generic compiler translates the validated description into SimPy storage containers, processing processes, rover processes, and a power-accounting process. Resource ports become inventories and transfers, routes become transport operations, and the selected equations determine the quantities, delays, and energy effects applied during each event. This common representation has been exercised with simple cargo-relay, ice-processing, and water-transfer scenarios in addition to the ISRU preset. It demonstrates that the architecture is no longer conceptually limited to LOX production and that a new operational scenario can be assembled from a SysML model, user-defined equations, and routes rather than from a new dedicated Python simulation script.

\begin{figure}[htbp]
    \centering
    \includegraphics[width=0.92\linewidth]{figures/generic_processor_behavior}
    \caption{Role-based generic scenario compiler: a Processor example. The identified class determines the standard operational behavior and the equations needed to represent it.}
    \label{fig:generic-processor-behavior}
\end{figure}

The equation builder is intentionally constrained. It uses a safe parser rather than arbitrary Python evaluation and accepts only supported mathematical operations and named variables. An equation affects the simulation only when its output maps to an implemented runtime behavior, such as a resource output, stored-resource level, processing time, travel time, energy consumption, power demand, or power generation. This prevents an equation shown in the interface from becoming an inactive annotation with no DES effect. The generic engine should nevertheless not yet be described as able to simulate every lunar mission. Complex multi-resource synchronization, richer storage constraints, stochastic operations, complete unit checking, and a broader library of operational behaviors remain open work.
